\section{辅助场方法中的重正化}\label{sec:renorm}

有效拉氏量式~(\ref{efflag}) 的重正化与带基础表示辅助场的情形完全类似~\cite{Bagan:1993zv,Ji:2017oey}，并且足以使基本场的格林函数有限。然而，在重正化式~(\ref{Eq:repl}) 中的 LCOs $J_1, {\overline J}_1$ 时会出现额外的复杂情况。
我们先在协变规范下考虑维数正规化（DR）中的微扰重正化。此时，规范不变的 LCOs 在重正化下一般会与具有相同或较低质量量纲的算符混合。混合算符可以是如下三种类型：1）规范不变算符；2）BRST 正合算符；
3）由运动方程（EOM）而为零的算符~\cite{Collins:1984xc}。对于由一个夸克和一个基础表示辅助场构成的 LCO，这种混合不会发生，因为那里出现的算符已具有最低量纲，而且如果手征对称性得以保持，同量纲算符之间的混合是被禁止的~\cite{Chen:2017mie}。但在当前情形这不再成立。

在 DR 中，由 BRST 对称性允许与 $J_1^{\mu\nu}$ 混合的算符是
\begin{align}\label{mixingop}
%J_1^{\mu\nu}&=F^{\mu\nu}_a\mathcal Q_a, \non\\
J_2^{\mu\nu}&=n_{\rho}(F^{\mu\rho}_a n^\nu-F^{\nu\rho}_a n^{\mu}){\mathcal Q}_a/n^2,\non\\
%J_3^{\mu\nu}&=(-in^\mu A_a^{\nu}+in^\nu A_a^\mu)(in\cdot D \mathcal Q)_a/n^2,
J_3^{\mu\nu}&=(-in^\mu A_a^{\nu}+in^\nu A_a^\mu)((in\cdot D-m) \mathcal Q)_a/n^2,
\end{align}
其中 $J_2^{\mu\nu}$ 规范不变；$J_3^{\mu\nu}$ 正比于 $\mathcal Q$ 的运动方程，因此在物理矩阵元中为零（当 $m=0$ 时上述混合模式已在文献~\cite{Dorn:1981wa} 中给出）。

存在混合时，重正化后的算符可写成三角形的混合形式
\begin{align}\label{oprenormmix}
\begin{pmatrix}
 J_{1, R}^{\mu\nu} \\  J_{2, R}^{\mu\nu} \\ J_{3, R}^{\mu\nu}
\end{pmatrix}
&=
\begin{pmatrix}
Z_{11} & Z_{12} & Z_{13}\\ 0 & Z_{22} & Z_{23} \\ 0 & 0 & Z_{33}
\end{pmatrix}
\begin{pmatrix}
 J_1^{\mu\nu}\\  J_2^{\mu\nu} \\ J_3^{\mu\nu}
\end{pmatrix}.
\end{align}
式~(\ref{oprenormmix}) 中的重正化常数并非全都独立。当 $\nu=z$ 时，$J_2^{z\mu}$ 与 $J_1^{z\mu}$ 简并，
因此二者必须具有相同的重正化，这意味着
\beq\label{Zrelation}
Z_{11}+Z_{12}=Z_{22}, \hspace{5em} Z_{13}=Z_{23}.
\eeq
这与明确的单圈计算一致~\cite{Dorn:1981wa}。
上述结果表明 $J_1^{\mu\nu}$ 的各分量独立重正化，相应的简化重正化方程为：
\begin{align}\label{oprenormmixsim2x2}
\begin{pmatrix}
 J_{1, R}^{z\mu} \\ J_{3, R}^{z\mu}
\end{pmatrix}
&=
\begin{pmatrix}
Z_{22} & Z_{13} \\  0 & Z_{33}
\end{pmatrix}
\begin{pmatrix}
 J_1^{z\mu} \\ J_3^{z\mu}
\end{pmatrix};\hspace{1em}
J_{1, R}^{ti}=Z_{11} J_1^{ti}.
\end{align}
$J_{1}^{ij}$ 与 $J_{1}^{ti}$ 具有相同的重正化。($J_{1}^{ti}$, $J_{1}^{ij}$) 这一组与 $J_1^{z\mu}$ 这一组算符之所以有不同的重正化，是因为四矢量 $n^\mu$ 的存在破坏了洛伦兹对称性。注意 $J_2^{\mu i}$ 并不独立，以下将其忽略。

我们的计算表明，$J_1^{\mu\nu}$ 的单圈修正中不出现线性发散。实际上，该算符的幂次发散结构已在文献~\cite{Wang:2017qyg,Wang:2017eel} 中用简单的截断正规化研究过。处理截断正规化中的幂次发散时必须谨慎，因为它一般会破坏 QCD 的规范不变性（格点截断除外），并可能掩盖理论中真实的幂次发散结构。为避免这一点，我们选择在 DR 中工作，并通过围绕 $d=3$ 展开来追踪线性发散——它们表现为 $d=3$ 处的极点。这样我们便能以规范不变的方式提取线性发散，并确认 $J_1^{\mu\nu}$ 的单圈修正中不存在线性发散。若采用规范不变的调节子，其实从文献~\cite{Wang:2017eel} 也能得到同样的结论。我们已经按照上述步骤明确验证了这一点。

在式~(\ref{efflag}) 中，我们还为辅助场 $\mathcal Q$ 引入了“剩余质量”项。其目的在于追踪可能的质量重正化：它在 DR 中围绕 $d=4$ 展开时也许不出现，但在诸如格点正规化之类的截断正规化中却会出现。在截断正规化中，辅助“重夸克”的有效质量即使在一阶不存在，也会由辐射修正产生。
这正是我们当前情形所发生的事实：我们有 $m=\delta m\sim{\cal O}(1/a)$，其中 $a$ 是截断或格距的倒数~\cite{Maiani:1991az}。在微扰论中，$\delta m$ 从 $O(\alpha_s)$ 阶开始出现。
如下文所示，这一质量项的作用是吸收积掉辅助场时威尔逊线自能产生的幂次发散。除此之外，理论中不存在其他幂次发散。此外，按量纲理由，也不存在能与前文讨论的 $J_i^{\mu\nu}$ 混合的其他反对称算符。因此在一个规范不变的截断方案中，算符重正化与 DR 中相同。

文献~\cite{Ji:2017oey} 曾证明，拉氏量的 BRST 不变性要求 $m$ 依赖于 $n$ 的号差~\cite{Dorn:1986dt}：对类光 $n^\mu$ 质量为零，对类时 $n^\mu$ 质量为实数。当 $n^\mu$ 为类空时，$m=\delta m=i\overline{\delta { m}}$ 为虚数。$\overline{\delta { m}}$ 在 $O(\alpha_s)$ 阶的值可以很容易地由文献~\cite{Ji:2017oey} 得到：$\overline{\delta { m}}={\pi\alpha_s C_A}/(2a)$。
